\scanpage{108}
\mspage{104}
% archive sheet number: 107

auxquels j'ai fait allusion ont \unsure{marché}, dans A), pour nous dire que la commutation des $\alpha_*$ à $|\unsure{\circ|0|\bullet}|$ implique les autres propriétés (sous les hyp. habituelles de \unsure{commensurabilité}). Donc l'interprétation géom. attachée à ``$\unsure{u_\alpha|u_*}$ commute à $|\circ|$'' sera bien celle d'un ``comorphisme'' $f^{-1}$, suivi d'une application injective. Quand \uline{de plus} $u_\alpha(1)=1$, on s'attendra donc à avoir commutation à $\unsure{C|\complement}$ (et on pourra le prouver, moyennant des conditions de \unsure{constructibilité} \unsure{convenables})

\begin{quote}
\noindent\vrule\quad\begin{minipage}{0.88\linewidth}
NB. La commutation à $C$ implique tjrs que $1\longrightarrow1$, ainsi que la compatibilité à $|\circ|$.
\end{minipage}
\end{quote}

\uuline{Question} Si $\alpha_*$ \uline{et} $\beta_*$ commutent à $C$, est-ce des iso. de spatios, \inserted{\unsure{(oui.)|(oui.}} (inverses l'un de l'autre) ?

Notons enfin que l'équivalence de (i) et (ii) dans (B) se vérifie trivialement dans tout spatios.
